\documentclass[12pt, a4paper]{article}

% --- 基础宏包设置 ---
\usepackage[utf8]{inputenc}
\usepackage{xeCJK}          % 处理中文字体（需使用 XeLaTeX 编译）
\usepackage{geometry}       % 页边距设置
\geometry{left=2.5cm, right=2.5cm, top=3cm, bottom=3cm}
\usepackage{cite}           % 引用文献管理
\usepackage{setspace}       % 行间距
\onehalfspacing             % 1.5倍行距
\usepackage{hyperref}       % 超链接
\hypersetup{colorlinks=true, linkcolor=blue, citecolor=blue, urlcolor=blue}
\usepackage{titlesec}       % 标题格式定制

% --- 标题信息 ---
\title{\textbf{基于卷积神经网络与Transformer融合架构的脑肿瘤多序列MRI影像分类研究}}
\author{}
\date{}

\begin{document}

\maketitle

% --- 引言部分开始 ---
\section{引言}

脑肿瘤作为中枢神经系统（CNS）中致死率最高的疾病之一，其精准的病理分类对于临床治疗方案的选择及预后评估具有至关重要的意义。根据 2021 年世界卫生组织（WHO）发布的第五版中枢神经系统肿瘤分类指南 \cite{Louis2021WHO}，脑肿瘤的诊断已从传统的形态学描述转向分子标志物与影像学特征深度结合的新阶段。磁共振成像（MRI）凭借其卓越的软组织对比度及多序列成像能力（如 T1-weighted, T2-weighted 及 FLAIR），能够清晰展示脑部的解剖结构与病理形态，已成为评估胶质瘤、脑膜瘤及垂体瘤等原发性肿瘤的首选非侵入性检查手段 \cite{Vaskovic2023Normal}。然而，脑肿瘤在影像学上表现出极高的异质性，且发病率的逐年攀升使得放射科医生面临沉重的判读压力，极易导致疲劳性的误诊或漏诊 \cite{Younis2023Survey}。因此，开发能够自动化、精准化分类脑肿瘤的深度学习模型，已成为当前人工智能与神经肿瘤学交叉领域的研究热点，旨在提升诊断的客观性与效率 \cite{Khalighi2024AI}。

放射科医生在进行 MRI 判读时，通常遵循一个从“局部病灶细节”到“全局空间关联”的综合推理过程。首先，医生需通过观察病灶的微观纹理和边缘特征来判断其性质。例如，恶性胶质瘤常表现出明显的浸润性生长特征，导致病灶边界模糊且形态不规则 \cite{Hakyemez2005Glioma}；而良性脑膜瘤则通常具有清晰的边界，并伴有特异性的“脑膜尾征”（Dural tail sign） \cite{Starr2017Meningioma}。卷积神经网络（CNN）凭借其局部感受野和权重共享的归纳偏置，在捕捉这些底层纹理特征和信号强度突变（Hyperintense/Hypointense）方面表现卓越。其次，肿瘤的生长往往伴随“占位效应”（Mass Effect），即巨大病灶会对周围正常组织产生挤压，导致侧脑室变形、中线偏移或严重的血管源性水肿 \cite{Smirniotopoulos2020Differential}。这种全局性的解剖结构扭曲是临床医生判断病变程度及占位性质的关键逻辑 \cite{Cleveland2024Brain}。

在具体的脑肿瘤分类任务中，早期的研究主要依赖于计算机辅助诊断（CAD）系统。这些系统通常利用手工设计的影像组学（Radiomics）特征，如灰度共生矩阵（GLCM）提取的纹理特征，并结合支持向量机（SVM）或随机森林等传统机器学习算法进行分类 \cite{Lohmann2020Radiomics}。然而，这类方法极度依赖专家的先验知识进行特征筛选，泛化能力有限。

随着深度学习技术的兴起，卷积神经网络（CNN）凭借其强大的特征自学习能力逐渐成为该领域的主流，多位研究者在 Sartaj 等公开数据集上利用 VGG、ResNet 及 DenseNet 等架构实现了显著的性能提升 \cite{Younis2023Survey}。研究重点随后转向对 CNN 架构的深度优化与鲁棒性增强：Devi \cite{Devi2021Detection} 利用网格搜索算法实现了超参数的自动化配置；Banerjee 等人 \cite{Banerjee2023Hyperparameter} 则通过系统的微调证明了深度卷积网络在识别肿瘤早期病灶中的潜力。为了进一步提升性能，Sharma \cite{Sharma2024Enhanced} 通过优化 Dropout 层与扁平化结构，显著增强了 CNN 处理高维影像时的泛化力；Sakthi 等人 \cite{Sakthi2023Deep} 提出的深度卷积框架进一步验证了端到端学习在区分不同风险等级肿瘤中的高效性。

为了解决医学影像标注成本高、样本量有限的挑战，近期研究开始广泛采用迁移学习与增强技术。Govind 等人 \cite{Govind2025Robust} 提出了一种基于迁移学习的鲁棒分类方法，利用预训练权重引导模型快速捕获复杂的病理特征。然而，传统卷积算子在建模解剖结构长程依赖上的局限性促使研究者转向视觉 Transformer（ViT）架构。Ding 等人 \cite{Ding2024VIT} 证明了 ViT 在极小样本数据集下依然能保持极高的分类精度，其自注意力机制能够敏锐捕捉区域间的语义关联。

尽管 CNN 在局部纹理提取上表现卓越，但其固有的局部感受野局限性使其难以捕捉解剖结构间的长程依赖。受大规模视觉预训练模型成功的启发，视觉 Transformer（ViT）开始被引入医学影像领域 \cite{Chen2021TransUNet}。初步研究表明，ViT 能够通过自注意力机制捕捉到 CNN 容易忽略的解剖结构间关联。Ding 等人 \cite{Ding2024VIT} 证明了 ViT 在小样本数据集下依然能保持极高的准确率，能够敏锐捕捉区域间的语义联系。

近来，ViT 的研究进一步向复杂场景与混合架构演进。针对临床中棘手的“多病灶”情况，Vani 等人 \cite{Vani2025Efficient} 利用 ViT 的长程建模能力解决了传统架构容易丢失全局空间细节的问题。为了增强特征表达，Hameed 等人 \cite{Hameed2025Feature} 创新性地将 ViT 与图卷积网络（GCN）结合，实现了对病灶空间拓扑结构的建模。此外，Ahmed 等人 \cite{Ahmed2024Hybrid} 提出的 ViT-GRU 混合模型不仅提升了时空特征提取效率，还引入了可解释性 AI 技术，为医疗决策提供了透明的证据支持。

尽管上述相关工作展示了卷积神经网络与视觉 Transformer 在脑肿瘤识别领域各自取得的显著成就，但在实际的临床辅助诊断场景中，单纯依赖单一架构模型仍面临严峻挑战。一方面，传统的 CNN 虽然在捕捉病灶边缘纹理等低级细节上表现卓越 \cite{Sharma2024Enhanced, Sakthi2023Deep}，但其固有的卷积核感受野限制了模型对病变组织与远端脑部解剖结构（如侧脑室压迫、中线偏移等“占位效应”）间长程依赖关系的建模能力。另一方面，单纯的 ViT 架构虽能通过全局自注意力机制建立跨区域的语义关联 \cite{Ding2024VIT}，但其往往缺乏归纳偏置，在处理医学影像这类对局部纹理高度敏感的小规模数据集时，收敛效率与细节表征能力仍有待提升 \cite{Khalighi2024AI}。

综上所述，脑肿瘤 MRI 的精准判读不仅需要对局部肿瘤浸润边缘进行细致观察，还需对由于“占位效应”引起的全局解剖结构扭曲进行空间推理。为了实现两种表征能力的互补，本研究设计并实现了一种串行式（CNN $\rightarrow$ Transformer）融合架构。该架构旨在利用 CNN 模块作为特征感官前端，发挥其在纹理提取上的感官优势；随后通过 Transformer 编码器对提取的特征序列进行全局依赖建模。本文后续将详细阐述该模型的设计逻辑、超参数优化过程及消融实验结果，旨在通过局部与全局特征的有机融合，为脑肿瘤的精准分类提供一种更具鲁棒性的技术方案。


\bibliographystyle{unsrt} % 或者是 unsrt (按引用顺序排序), IEEEtran 等
\bibliography{bi}

\end{document}